\documentclass{article}
\usepackage{amsmath, amssymb, booktabs, xcolor}

\begin{document}

\section*{Pushing--Medium: Cross-Force Interactions and Observational Hierarchy}

\textbf{Status Tags:}
\textcolor{blue}{[CALIBRATED]}, \textcolor{green}{[GR-MAPPED]}, \textcolor{orange}{[DIMENSIONAL]}, \textcolor{red}{[CONSTRAINED]}, \textcolor{gray}{[RULED OUT]}

\section*{1. What PM Allows That GR Does Not Include}

\subsection*{1.1 Structural Differences}

In General Relativity, gravity is geometric (spacetime curvature) and
electromagnetism is a gauge field on that spacetime. The only coupling between
them is via the stress--energy tensor $T_{\mu\nu}$.

In Pushing--Medium, both are excitations of the same medium ($\phi$ and
$\mathbf u$ fields). This structure allows \emph{field-level couplings} that do
not appear in GR. Any analogous GR effect would have to arise via
$T_{\mu\nu}$—which is typically negligible.

\subsection*{1.2 Four Interaction Mechanisms}

\subsubsection*{A. Unified Index Superposition}

\textbf{PM structure:}



\[
\ln n_{\text{tot}} = \mu_G\,\Phi_G + \mu_E\,\Phi_E - \alpha_{\text{plasma}}\,n_e
\]



\textbf{Consequence:} Mass, charge, and plasma all contribute to the same
refractive index field. A charged mass produces both gravitational and
electrostatic lensing contributions.

\textbf{GR:} Only total mass-energy lenses (via metric curvature). Charge
contributes via $E^2/(2\varepsilon_0)$ energy density, which is negligible for
macroscopic objects.

\textbf{Status:}

\begin{itemize}
    \item $\mu_G$: \textcolor{blue}{[CALIBRATED]} --- $1.486 \times 10^{-27}$ m$^2$/kg from weak-field tests.
    \item $\mu_E$: \textcolor{orange}{[DIMENSIONAL]} --- Estimated $\sim 10^{-37}$ m$^2$/C; unconstrained by data.
    \item $\alpha_{\text{plasma}}$: \textcolor{orange}{[DIMENSIONAL]} --- Estimated $\sim 10^{-15}$ m$^3$ from standard plasma optics; choice to include in gravity sector is \emph{modeling ansatz}.
\end{itemize}

\subsubsection*{B. Current-Induced Dragging}

\textbf{PM structure:}



\[
\nabla^2 \mathbf u = -A_M\,\mathbf J_M - A_q\,\mathbf J_q
\]



\textbf{Consequence:} Both mass currents ($\mathbf J_M = \rho_M \mathbf v$) and
charge currents ($\mathbf J_q$) source flow fields that drag light via:



\[
\frac{d\mathbf r}{ds} = c\,\hat{\mathbf k} + \mathbf u_{\text{tot}}.
\]



\textbf{GR:} Frame-dragging comes from rotating mass-energy (metric perturbation
$g_{t\phi}$). Electric currents contribute via $T^{0i}$ (stress-energy flux),
which for lab fields is $\sim 10^{-40}$ smaller than PM's $A_M$ term.

\textbf{Status:}

\begin{itemize}
    \item $A_M$: \textcolor{green}{[GR-MAPPED]} --- $G/c^2 \approx 7.4 \times 10^{-28}$ m$^3$/(kg$\cdot$s) from matching Lense--Thirring formula.
    \item $A_q$: \textcolor{red}{[CONSTRAINED]} --- Lab magnet null results require $A_q < 10^{-10}$ m$^3$/(C$\cdot$s). See Section 3 for derivation.
    \item $\kappa_B$ (magnetic potential coupling): \textcolor{gray}{[RULED OUT]} --- Would predict $\alpha \sim 10$ rad near lab solenoids. Set to zero.
\end{itemize}

\subsubsection*{C. Plasma in the Gravitational Index}

\textbf{PM structure:} Electron density $n_e$ contributes to the refractive
index alongside gravitational and electrostatic terms.

\textbf{GR:} Plasma affects light via electromagnetic dispersion
($n_{\text{plasma}}(\omega)$), not via metric geometry. Gravitational lensing
and plasma dispersion are independent.

\textbf{Status:} \textcolor{orange}{[MODELING ANSATZ]} --- Including plasma in
$\ln n_{\text{tot}}$ is a PM design choice. Could lead to frequency-dependent
lensing corrections; not yet tested.

\subsubsection*{D. Dynamic Field Mixing (Advection)}

\textbf{PM structure:}



\[
\frac{\partial \phi}{\partial t} + \nabla\cdot(\phi\,\mathbf u) = S_\phi + \ldots
\]



Scalar field $\phi$ and flow field $\mathbf u$ are dynamically coupled. Since
$\mathbf u$ includes contributions from multiple sources, energy can transfer
between gravitational and EM excitations via the $-\phi\,\nabla\cdot(\phi\,\mathbf u)$ interaction.

\textbf{GR:} Gravitational field $g_{\mu\nu}$ and EM field $F_{\mu\nu}$ evolve
independently except via one-way coupling through $T_{\mu\nu}$.

\textbf{Status:} \textcolor{orange}{[PREDICTED BY STRUCTURE]} --- Required by PM
Lagrangian; not yet tested in dynamic regimes.

\section*{2. The Observational Hierarchy}

\begin{center}
\begin{tabular}{llllp{4cm}}
\toprule
\textbf{Symbol} & \textbf{Sector} & \textbf{Value/Bound} & \textbf{Status} & \textbf{Determined By} \\
\midrule
$k_{\text{Fizeau}}$ & Kinematic & $1.0$ & \textcolor{blue}{[CALIBRATED]} & Moving lens tests \\
$k_{\text{TT}}$ & Kinematic & $1.0$ & \textcolor{blue}{[CALIBRATED]} & Time delay tests \\
\midrule
$\mu_G$ & Gravity & $1.486 \times 10^{-27}$ & \textcolor{blue}{[CALIBRATED]} & Weak-field deflection \\
$A_M$ & Gravity & $7.426 \times 10^{-28}$ & \textcolor{green}{[GR-MAPPED]} & Lense--Thirring match \\
\midrule
$\mu_E$ & EM & $\sim 10^{-37}$ & \textcolor{orange}{[DIMENSIONAL]} & Unconstrained \\
$A_q$ & EM & $< 10^{-10}$ & \textcolor{red}{[CONSTRAINED]} & Lab magnet null \\
$\alpha_{\text{plasma}}$ & Plasma & $\sim 10^{-15}$ & \textcolor{orange}{[DIMENSIONAL]} & Standard plasma \\
$\kappa_B$ & EM & $\approx 0$ & \textcolor{gray}{[RULED OUT]} & Lab magnet null \\
\midrule
$c_\phi$ & Dynamics & $c$ & \textcolor{orange}{[ASSUMED]} & Pending GW tests \\
$c_u$ & Dynamics & $c$ & \textcolor{orange}{[ASSUMED]} & Pending GW tests \\
\bottomrule
\end{tabular}
\end{center}

\textbf{Units:}
$\mu_G$ (m$^2$/kg), $A_M$ (m$^3$/(kg$\cdot$s)), $\mu_E$ (m$^2$/C), $A_q$
(m$^3$/(C$\cdot$s)), $\alpha_{\text{plasma}}$ (m$^3$), $\kappa_B$
((m/s)/(T$\cdot$m$^2$)).

\subsection*{2.1 The Coupling Hierarchy}

Given the PM structure, the coupling constants are \emph{free parameters}.
Observations impose:



\[
\frac{A_q}{A_M} < 4 \times 10^{17}
\quad\text{(charge currents couple $\gtrsim 10^{17}\times$ more weakly)}.
\]



This hierarchy is an \emph{empirical property of the medium}, not a theoretical
necessity. It explains why gravitational effects dominate in all tested regimes
and why GR (which omits EM-gravity cross-terms) succeeds.

\section*{3. Derivation of Observational Bounds}

\subsection*{3.1 Constraint on $A_q$: Laboratory Solenoid Null Result}

\textbf{Setup:}

\begin{itemize}
    \item Superconducting solenoid: $I = 10^6$ A, length $L = 1$ m.
    \item Light beam passes at distance $r = 0.1$ m.
    \item Path length through interaction region: $\ell_{\text{path}} = 1$ m.
    \item Detection threshold: $\alpha_{\text{det}} \sim 10^{-10}$ rad (optical interferometry limit).
\end{itemize}

\textbf{PM prediction:}

If charge currents couple via $\nabla^2\mathbf u_q = -A_q\,\mathbf J_q$, the
flow field magnitude is:



\[
|\mathbf u_q| \sim A_q\,I\,\frac{L}{r}.
\]



Angular deflection:



\[
\alpha \sim \frac{|\mathbf u_q|}{c}\,\frac{\ell_{\text{path}}}{r}
= A_q\,I\,\frac{L\,\ell_{\text{path}}}{c\,r^2}.
\]



\textbf{Observational constraint:}

No deflection observed above $\sim 10^{-10}$ rad. Therefore:



\[
A_q < \frac{\alpha_{\text{det}}\,c\,r^2}{I\,L\,\ell_{\text{path}}}
\sim \frac{10^{-10} \times 3 \times 10^8 \times 0.01}{10^6 \times 1 \times 1}
\sim 3 \times 10^{-10}\,\text{m}^3/(\text{C}\cdot\text{s}).
\]



\textbf{Comparison to gravity:}



\[
\frac{A_q}{A_M} < 4 \times 10^{17}.
\]



\textbf{Interpretation:} This is an \emph{order-of-magnitude illustrative
bound} based on reasonable lab parameters. Tighter bounds could be derived from:

\begin{itemize}
    \item MRI systems (routine operation with no anomalous optical effects),
    \item tokamak diagnostics (optical systems in mega-ampere environments),
    \item precision magnetometry (interferometric setups near strong fields).
\end{itemize}

For practical PM implementations: $A_q \approx 0$.

\subsection*{3.2 Constraint on $\kappa_B$: Magnetic Potential Coupling}

\textbf{Initial formulation:} $\mathbf u_B = \kappa_B\,\mathbf A$.

\textbf{Problem:} For the same solenoid, $|\mathbf A| \sim \mu_0 I L / (4\pi r)
\sim 10^{-2}$ T$\cdot$m. With $\kappa_B \sim c$:



\[
|\mathbf u_B| \sim c \times 10^{-2} \sim 10^6\,\text{m/s}.
\]



Predicted deflection: $\alpha \sim 10$ rad.

\textbf{Observational verdict:} Not seen.

\textbf{Resolution:} Set $\kappa_B = 0$. The magnetic vector potential does
\emph{not} couple directly to the PM flow field. This avoids double-counting
(since $\mathbf A$ is already a functional of $\mathbf J_q$) and respects
observations.

\section*{4. Structural Resolution: Sources, Not Potentials}

The correct PM coupling is:



\[
\mathbf S_u = A_M\,\mathbf J_M + A_q\,\mathbf J_q.
\]



\textbf{Rationale:}

\begin{itemize}
    \item Couples at \emph{source level} (current densities), not field level (potentials).
    \item Treats mass and charge currents on equal structural footing.
    \item Avoids double-counting ($\mathbf A$ is derived from $\mathbf J_q$).
    \item Respects ``no special cases'' principle.
    \item Consistent with all laboratory null results.
\end{itemize}

\textbf{Comparison to electromagnetism:}

In Maxwell's equations, sources ($\rho$, $\mathbf J$) determine fields
($\mathbf E$, $\mathbf B$). We don't couple sources to \emph{both} $\mathbf J$
and derived fields like $\mathbf B$. PM follows the same principle.

\section*{5. What Is Solid, What Is Aspirational}

\subsection*{5.1 Solid (Observationally Anchored)}

\begin{itemize}
    \item PM gravity reproduces GR weak-field predictions exactly.
    \item $\mu_G$, $k_{\text{Fizeau}}$, $k_{\text{TT}}$ are calibrated to data.
    \item $A_M$ matches GR frame-dragging structure.
    \item $\kappa_B$ is ruled out by lab null results.
    \item $A_q$ has an observational upper bound (though still $\sim 27$ orders above dimensional estimate).
\end{itemize}

\subsection*{5.2 Aspirational (Structural Predictions)}

\begin{itemize}
    \item $\mu_E$ coupling exists in structure but is unconstrained.
    \item Charge-dependent lensing predicted but not tested.
    \item $\alpha_{\text{plasma}}$ term is modeling ansatz (putting plasma in gravity sector).
    \item Dynamic field mixing ($\nabla\cdot(\phi\,\mathbf u)$) predicted but not tested.
    \item Multiple wave speeds ($c_\phi$, $c_u$) assumed equal to $c$ pending GW data.
\end{itemize}

\subsection*{5.3 What PM Does NOT Claim}

\begin{itemize}
    \item PM does \emph{not} claim to have measured $\mu_E$ or $A_q$ directly.
    \item PM does \emph{not} claim EM-gravity cross-terms are large or easily observable.
    \item PM does \emph{not} claim to explain anomalies unexplained by GR (yet).
\end{itemize}

What PM \emph{does} claim:

\begin{itemize}
    \item The unified structure is mathematically consistent.
    \item It reproduces GR where tested.
    \item It makes testable predictions at precision frontiers.
    \item It respects ``no special cases''—same equations, different couplings.
\end{itemize}

\section*{6. The Observational Message}

The hierarchy $A_M \gg A_q$ (by $\gtrsim 10^{17}$) tells us:

\begin{quote}
\emph{The PM medium responds vastly more strongly to mass than to charge.}
\end{quote}

This is not fine-tuning. Given the PM structure, the couplings are free
parameters—observations then \emph{measure} them and impose the hierarchy.

Analogy: The Standard Model has three gauge couplings ($g_1$, $g_2$, $g_3$)
with a hierarchy $\alpha_{\text{strong}} \gg \alpha_{\text{EM}}$. This is
empirical, not derived. PM's $A_M \gg A_q$ is the same kind of fact.

\section*{7. Why This Is Not a Problem}

The small/zero values for EM couplings mean:

\begin{itemize}
    \item Gravitational sector dominates (as observed).
    \item PM reduces to ``pure gravity'' in all tested regimes.
    \item Cross-terms exist \emph{in principle} but are negligible \emph{in practice}.
    \item GR's simpler structure (omitting cross-terms) is empirically justified.
\end{itemize}

But PM remains \emph{structurally richer}:

\begin{itemize}
    \item Predicts where to look for deviations (charged beams, extreme plasmas).
    \item Offers unified framework for future precision tests.
    \item Shows how forces \emph{could} be unified via medium excitations.
\end{itemize}

\section*{8. The Landmine: Documented and Resolved}

\subsection*{8.1 What Went Wrong Initially}

Naive dimensional analysis suggested $\mathbf u_B = \kappa_B\,\mathbf A$ with
$\kappa_B \sim c$. This predicts:



\[
|\mathbf u_B| \sim 10^6\,\text{m/s} \quad\Rightarrow\quad \alpha \sim 10\,\text{rad}
\]



near 1 MA lab solenoids.

\textbf{Not observed.}

\subsection*{8.2 The Fix}

\begin{itemize}
    \item \textbf{Drop:} $\kappa_B\,\mathbf A$ term (set $\kappa_B = 0$).
    \item \textbf{Keep:} $A_M\,\mathbf J_M + A_q\,\mathbf J_q$ (source coupling only).
    \item \textbf{Reason:} $\mathbf A$ is derived from $\mathbf J_q$ via $\nabla^2\mathbf A = -\mu_0\mathbf J_q$. Coupling to both is double-counting.
\end{itemize}

\subsection*{8.3 Why This Is Better}

\begin{enumerate}
    \item Structurally cleaner (source-level coupling).
    \item Observationally consistent (no huge lab predictions).
    \item Philosophically correct (``no special cases''—all currents treated uniformly).
    \item Provides testable constraint on $A_q$.
\end{enumerate}

This resolution \emph{strengthens} PM by removing an ad-hoc, problematic term
and leaving a cleaner, more symmetric structure.

\section*{9. Honest Summary}

\subsection*{9.1 What PM Has Achieved}

\begin{itemize}
    \item A gravity theory that passes all GR tests in weak/semistrong regimes.
    \item Calibrated constants: $\mu_G$, $k_{\text{Fizeau}}$, $k_{\text{TT}}$.
    \item A unified structure that \emph{allows} cross-force interactions.
    \item Observational constraints showing these interactions are tiny.
\end{itemize}

\subsection*{9.2 What PM Has Not Yet Done}

\begin{itemize}
    \item Measured $\mu_E$, $A_q$, or $\alpha_{\text{plasma}}$ directly.
    \item Explained an anomaly that GR cannot explain.
    \item Tested in strong-field or dynamic regimes where PM and GR must differ.
\end{itemize}

\subsection*{9.3 What PM Offers}

\begin{itemize}
    \item A mathematically consistent alternative to GR.
    \item Structural unity (all forces as medium excitations).
    \item Testable predictions at precision frontiers.
    \item A framework for exploring gravity-EM unification.
\end{itemize}

\section*{10. Recommended Language for PM Presentations}

\subsection*{10.1 Avoid}

\begin{itemize}
    \item ``PM reveals interactions that cannot exist in GR'' $\rightarrow$ too absolute.
    \item ``This explains [unsolved problem]'' $\rightarrow$ unless you've actually fit the data.
    \item ``The hierarchy is natural'' $\rightarrow$ implies derivation; it's actually empirical.
\end{itemize}

\subsection*{10.2 Use}

\begin{itemize}
    \item ``PM allows field-level couplings not included in GR; analogous GR effects via $T_{\mu\nu}$ are negligible.''
    \item ``Observations constrain EM couplings to be $\gtrsim 10^{17}$ weaker than gravity couplings.''
    \item ``Given the PM structure, the couplings are free parameters; data then impose the hierarchy.''
    \item ``PM predicts [X]; testing this requires [Y precision/environment].''
\end{itemize}

\section*{11. When Someone Asks: ``Why doesn't this blow up in an MRI?''}

\textbf{Short answer:}

``The magnetic potential coupling ($\kappa_B\,\mathbf A$) is ruled out by
exactly that consideration—we set $\kappa_B = 0$. The remaining charge-current
coupling ($A_q\,\mathbf J_q$) is constrained by the same null results to be
$< 10^{-10}$ m$^3$/(C$\cdot$s), making it negligible.''

\textbf{Long answer:} See Section 3.

\section*{12. Future Work}

\subsection*{12.1 Immediate (Refine Bounds)}

\begin{itemize}
    \item Systematic review of lab magnet environments for tighter $A_q$ bound.
    \item Charged-particle beam deflection experiments to constrain $\mu_E$.
    \item Frequency-dependent cluster lensing to test $\alpha_{\text{plasma}}$ ansatz.
\end{itemize}

\subsection*{12.2 Medium-Term (Explore Predictions)}

\begin{itemize}
    \item Strong-field regime (where PM has no singularities, GR has horizons).
    \item Dynamic advection effects (where $\nabla\cdot(\phi\,\mathbf u)$ matters).
    \item Multi-messenger GW tests (for $c_\phi \neq c_u$ possibility).
\end{itemize}

\subsection*{12.3 Long-Term (Theoretical)}

\begin{itemize}
    \item Derive coupling ratios from PM microphysics (if any).
    \item Explore quantum corrections to couplings.
    \item Connect to other unification frameworks.
\end{itemize}

\section*{13. Bottom Line}

PM's unified field structure is:

\begin{itemize}
    \item \textbf{Mathematically consistent} (passes internal checks).
    \item \textbf{Observationally consistent} (no glaring contradictions).
    \item \textbf{Conceptually unified} (``no special cases'').
    \item \textbf{Empirically constrained} (hierarchy imposed by data).
\end{itemize}

The cross-force interactions exist in the formalism but are naturally suppressed
by the coupling hierarchy. This is \emph{not a bug but a feature}—it shows PM is
richer than GR while remaining empirically viable.

The path forward: measure the unconstrained couplings and test the structural
predictions at the precision frontier.

\end{document}
